\section{Exact arithmetic checks}\label{app:checks}

\subsection{Logarithm bounds and endpoint comparisons}
\label{app:endpoint-checks}

Here is the finite rational calculation used in
Section~\ref{sec:count}. For a rational $x\ge1$, write
$x=2^j y$ with $j\ge0$ and $1\le y<2$, and put
$z=(y-1)/(y+1)$. With $M=18$, define
\[
 L_M(y)=2\sum_{h=0}^{M-1}\frac{z^{2h+1}}{2h+1},\qquad
 R_M(y)=\frac{2z^{2M+1}}{(2M+1)(1-z^2)}.
\]
The positive series $\log y=2\operatorname{atanh}z$ gives
\[
 L_M(y)\le\log y\le L_M(y)+R_M(y).
\]
Applying the same formula to $2$ gives bounds for $\log2$.
Adding $j$ times those bounds gives rational numbers
$\underline{\log}(x)$ and $\overline{\log}(x)$ enclosing
$\log x$. The remainder follows by replacing all denominators
in the tail by $2M+1$ and summing the resulting geometric series.

The checker computes $\varphi(d)$ by counting coprime integers and
then forms $D_k$. For each closed endpoint $s_k$, it computes
\[
 2\left\lceil D_k\left(33+\frac{45}{2}
                     \underline{\log}(k)\right)\right\rceil,
 \qquad
 2\left\lceil D_k\left(33+\frac{45}{2}
                     \overline{\log}(k)\right)\right\rceil.
\]
These integers agree in every required case, thereby isolating the
endpoint exactly. It then performs the following rational comparisons:
\begin{center}
\begin{tabular}{@{}lrr@{}}
\toprule
Comparison&Number of cases&Index of smallest margin\\
\midrule
Unrounded, $26\le k\le255$&230&29\\
Rounded, $3\le k\le25$&23&3\\
\bottomrule
\end{tabular}
\end{center}
After subtraction of $968/6615$, the margins are strictly greater
than $24529/10^8$ and $14325/10^8$, respectively, as used in the
proof. All calculations use integer or rational arithmetic. The
checker also verifies the constants in
\eqref{eq:tail-T}--\eqref{eq:tail-G}, the two small floors at
$s=160,364$, and the positive coefficient bound at the end of
Section~\ref{sec:count}.

The elementary bounds on $\pi$ can be verified independently of
decimal approximations. For the upper bound,
\[
 0<\int_0^1\frac{x^4(1-x)^4}{1+x^2}\,dx=\frac{22}{7}-\pi.
\]
For the lower bound $\pi>157/50$, the checker uses Machin's identity
$\pi=16\arctan(1/5)-4\arctan(1/239)$ with eight terms of the
alternating arctangent series and its next-term error bound.

\subsection{Checks of the construction identities}

The positive-definiteness proof in Section~\ref{sec:gram} works for
every $B$. We also check its normalization and the arithmetic
identities by separate finite calculations. At each of the ten
parameter choices in Section~\ref{sec:endpoints}, two partial-fraction
algorithms over the field with $1\,000\,003$ elements agree with
an exact rational calculation of the Gram matrix.

The rational calculation uses explicit moments. Every element of
$\mathcal A$ occurs in $\Lambda$, because the kernel poles are
numerator zeros. With all fractions reduced, put
\[
 \Delta=\prod_{\alpha=v/d\in\Lambda}d^2,\qquad
 V(x)=\prod_{\alpha=v/d\in\Lambda\setminus\mathcal A}(d^2x+v^2).
\]
Then $\prod_{\alpha\in\Lambda}(x+\alpha^2)/Q(x)=V(x)/\Delta$.
For $0\le k<S$, the elementary beta integral gives
\begin{equation}\label{eq:moments}
 \frac1{2\pi}\int_{-\infty}^{\infty}
                   \frac{y^{2k}}{(y^2+4)^S}\,dy
 =2^{2k-4S+2}\frac{(2k)!(2S-2k-2)!}{k!(S-k-1)!(S-1)!}.
\end{equation}
Thus $M_{b,j}/C_1$ is $\sigma/\Delta$ times the moment functional
in \eqref{eq:moments} applied to $xV(x)h_b(x)h_j(x)$. The largest
moment index is $DL/2+N-1<S$, so every integral used is convergent.

At the ten endpoints, the computations check $385$ rational Gram
entries, $55$ positive leading principal minors and ten integer
adjugate identities. Full rational computations at
$(B,s,r)=(1,20,3),(2,24,3)$ and $n=1,17,19$ also check the
coefficient and constant congruences and grouped integrality.
These smaller parameter sets test identities; decay is checked
at the stated theorem endpoints.

Run \path{checks/verify_construction.py} for the construction
identities and \path{checks/verify_endpoints.py} for the endpoint
comparisons. Both use only the Python standard library and should
be run without \texttt{-O}.

\subsection{Independent audit}

An independent mathematical audit reported on 5 November 2026
found no substantive error or gap in the 19-page revision identified
in the accompanying audit record. The reviewer checked the general
arguments and independently reproduced the finite calculations,
without running or importing the supplied checkers. The audit
recommended the explanation of uniformity now included in
Section~\ref{sec:profile-sum}; no mathematical correction was requested.
It was not formal verification or a fresh audit of all cited literature.
